\section{POLÍTICAS INSTITUCIONAIS NO ÂMBITO DO CURSO}

O \NOMECURSO do IFCE é um curso que vem ganhando renome no âmbito do estado, reconhecido pelos órgãos governamentais de ensino e pesquisa, pelo setor empresarial e, enfatize-se, pela sociedade. O curso também é bastante conhecido pelas pessoas que não estão relacionadas diretamente com a educação, mas que possuem consciência da importância dos bons cursos para a sociedade. Esse reconhecimento da sociedade é resultado de diversas políticas adotadas pelos dirigentes desde sua criação, das quais pode-se citar:

\begin{itemize}
	\item Contratação de professores em regime de 40 horas e dedicação exclusiva (DE) em meados da década de 2010, fundamentado na crença de que bons cursos só podem ser desenvolvidos com professores inteiramente dedicados a ele e em constante reciclagem;
	
	\item Institucionalização dos CEFETs, o que alavancou a pesquisa e a extensão nos atuais IFs;
	
	\item Política de incentivo a qualificação docente através de cursos a nível de pós-graduação (mestrado, doutorado e pós-doutorado) dos professores com afastamento temporário, podendo os mesmos realizarem seus estudos no Brasil ou no exterior, de forma a prover recursos humanos para a condução de pesquisas de ponta;
	
	\item Incentivo constante para a interação universidade/empresa, com o objetivo de transferir para a sociedade os conhecimentos adquiridos nas atividades de pesquisa bem como direcionar estas mesmas pesquisas à resolução de demandas reais, além estreitar laços que melhor permitam a inserção de nossos alunos no mercado de trabalho;
	
	\item Preocupação em sempre admitir profissionais de qualidade acadêmica em seu quadro.
\end{itemize}

O processo de pós-graduação dos professores está com cerca de 70\% concluído, e espera-se que até o final de 2019 esteja com 100\% de professores doutores (ver seção \ref{CORPODOCENTE}).

Os esforços também passam por melhorias e aumento da infraestrutura de ensino e pesquisa. Os laboratórios, (ver seção \ref{INFRA}, atualmente em número de 13), foram criados, ampliados e especializados e a interação com o setor empresarial vem se consolidando. A Tab. \ref{tab11} representa o número de projetos de pesquisa e de extensão aprovados pelos professores do eixo da indústria. 

\begin{longtable}{|p{1.5cm}|p{4.5cm}|p{4.5cm}|}
	\caption{Número de projetos de pesquisa e de extensão}    \label{tab11} 
	\vspace{-0.25cm}
	%\multicolumn{3}{c}{\textbf{Tabela 9. Número de projetos de pesquisa e de extensão}}    \label{tab11} 
	\\ \hline 
	\cellcolor{\cinza} Ano & \cellcolor{\cinza}Projetos de Pesquisa & \cellcolor{\cinza}Projetos de Extensão  \\ 
	\hline
	2013 &  9  &  15 \\ \hline
	2014 &  13  &  34 \\ \hline
	2015 &  07  &  18 \\ \hline
	2016 &  12  &  07 \\ \hline
	2017 &  19  &  03 \\ \hline
	Total &  59  & 77 \\ \hline
\end{longtable}

O ensino e a pesquisa, como deve a ser a regra, andam lado a lado no \NOMECURSO, propiciando o nascimento de um forte curso, uma vez que seus docentes, em função de suas atividades em pesquisa, sempre procuraram ministrar aos alunos, sólidos fundamentos da área da disciplina. Os laboratórios foram aparelhados com equipamentos diversos. Projetos em sua maioria de própria autoria da Gestão do Campus, mais também provindos de financiamento de empresas, do CNPq, do Finep e outros, forneceram os recursos para a aquisição dos diversos equipamentos hoje disponíveis, para a manutenção desses equipamentos, para a montagem de estruturas computacionais, bem como para a aquisição de material de consumo. Estes equipamentos estão, em sua maioria, listados na seção \ref{INFRA}. Ao mesmo tempo em que o apoio governamental direto, através do Ministério da Educação se restringia aos salários dos professores e a infraestrutura básica, os recursos dos projetos foram permitindo a contratação de pessoal técnico e também o pagamento de bolsas de iniciação científica, em muitos casos. Criou-se, assim, uma grande estrutura de pesquisa e de extensão, com reflexos diretos no ensino de graduação, através da iniciação científica e também pela própria participação do professor em sala de aula. 

A argumentação inicial de que a localização geográfica do IFCE Campus Maracanaú, em uma cidade praticamente provida de indústrias, seria benéfico ao curso, vem sendo comprovada no dia a dia do Campus, pois a influência do centro de pesquisa e de formação aqui instalado, foi muito além da fronteira municipal, ultrapassando a fronteira estadual e até nacional. Tal nacionalização e posterior internacionalização também vem sendo complementada pelo recém-criado Mestrado em Energias Renováveis do IFCE Campus Maracanaú, proposto e aprovado pelos professores do eixo da indústria, aqui mencionados, e do eixo da química e meio ambiente também presente no Campus. Não se tem empreendido um esforço direto, até o presente momento, para quantificar o tamanho desta influência na economia estadual e também nacional; entretanto sabe-se da sua importância nos meios acadêmicos e industriais. 

Assim, a qualidade de nossos alunos vem melhorando ano a ano, principalmente, após a intensificação de participação destes alunos como bolsistas de iniciação científica nos laboratórios. Esta atividade, além de fornecer subsídios de aprendizagem de conteúdos técnicos, extrapolando a sala de aula, e mostrando aplicações para os fundamentos lá vistos, permite ao aluno um aprendizado gradual e não aparente da cultura da pesquisa, do comportamento do engenheiro perante os problemas e acelera seu processo de maturidade, pois em geral ele convive com mestrandos e professores. Vale ressaltar que a inserção dos alunos de graduação nos laboratórios de ensino e pesquisa, permite o seu contato com estas atividades e com as atividades de extensão. Esses alunos são orientados e estão em contato direto com os alunos de mestrado, auxiliando-os a desenvolver os temas de pesquisa destes últimos. Os trabalhos de extensão, também são resolvidos nesses ambientes, com a participação de todos. A discussões que surgem, os experimentos que são montados, os resultados que são obtidos, criam uma boa atmosfera para o crescimento dos estudantes. Além de seu grupo de colegas da mesma fase, o aluno inserido no laboratório passa a fazer parte de outro grupo, dando-lhe mais maturidade, tendo como retorno o melhoramento do seu desempenho nas diversas disciplinas.

Outras atividades incentivadas ao longo do curso e que também motivam os alunos, são: Projetos AeroDesign e Enactus; Empresa Júnior; Centro Acadêmico; organização de eventos técnicos, culturais e esportivos; Cooperação Internacional; viagens de estudos; estágios; etc.

A Tabe. \ref{tab10} apresenta os investimentos financeiros relacionados ao \NOMECURSO desde o ano de 2012. 

\begin{longtable}{|p{1.0cm}|p{6.0cm}|p{3.0cm}|p{3.0cm}|}
	
	\caption{Investimentos financeiros relacionados ao curso}\label{tab10} 
	\vspace{-0.25cm}
	
	%\multicolumn{4}{c}{\textbf{Tabela 8. Investimentos financeiros relacionados ao curso}}\label{tab10} 
	\\ \hline 
	\cellcolor{\cinza}Ano & \cellcolor{\cinza}Equipamento & \cellcolor{\cinza}Empenho, Pregão, Contrato & \cellcolor{\cinza}Investimento [R\$]  \\ 
	\hline
	2012 &  Bloco da Oficina Mecânica & Nº40/2012  &  3.743.946,00 \\ \hline
	2013 &  Equipamentos para os Laboratórios  & PE 14/2013 e PE 06/2013 &  70.167,70 \\ \hline
	2014 &  Maquinário da Oficina Mecânica & PE 08/2014 &  938.260,88 \\ \hline
	2016 &  Maquinário da Oficina Mecânica  & PE 11/2016 & 135.004,76 \\ \hline
	2016 &  Livros & 2016NE800268, 2016NE800331, 2016NE800269 &  113.700,30 \\ \hline
	&  Total & & 5.001.079,64 \\ \hline
\end{longtable}

% APOIO AO DISCENTE ---------------------------------------------
